% TikZ 图：变式训练中的四分之一圆轨迹
% Day 22：洛伦兹力
\begin{tikzpicture}[scale=1.02,>=Latex,line cap=round,line join=round]
  \path[use as bounding box] (-1.65,-1.55) rectangle (2.80,1.45);
  \draw[thick,fill=blue!4] (0,-1.35) rectangle (1.65,1.25);
  \foreach \x in {0.38,0.92,1.46} {
    \foreach \y in {-0.85,-0.25,0.35,0.95} {
      \draw[blue!70!black,thick] (\x-0.05,\y-0.05) -- (\x+0.05,\y+0.05);
      \draw[blue!70!black,thick] (\x-0.05,\y+0.05) -- (\x+0.05,\y-0.05);
    }
  }
  \node[blue!70!black,font=\small] at (2.30,0.62) {$B$ 向里};

  \coordinate (O) at (0,0);
  \coordinate (A) at (0,-1);
  \coordinate (C) at (1,0);
  \draw[->,red,very thick] (-1.25,-1) -- (A);
  \node[red] at (-0.72,-1.20) {$v_1$};
  \draw[red,very thick,->] (A) arc (-90:0:1);
  \draw[->,red,very thick] (C) -- ++(0,0.88);
  \node[red,right] at (1.10,0.78) {$v_2$};

  \draw[gray,dashed,thick] (O) -- (A);
  \draw[gray,dashed,thick] (O) -- (C);
  \fill (O) circle (1.5pt) node[above left] {$O$};
  \fill[red] (A) circle (2pt);
  \node[red] at (-0.12,-1.30) {$A$};
  \node[red] at (0.34,-1.22) {$+q$};
  \fill[red] (C) circle (2pt) node[below right=2pt] {$C$};
  \draw[orange,thick,->] (0,-0.45) arc (-90:0:0.45);
  \node[orange] at (0.35,-0.34) {$90^\circ$};
\end{tikzpicture}
